\documentclass[11pt,oneside,a4paper]{book}
\usepackage[utf8]{inputenc}
\usepackage[spanish, es-tabla, es-lcroman]{babel}  
\usepackage{calc}
\usepackage{enumerate}    
\usepackage{etoolbox}
\usepackage{hyperref}                       
\usepackage{amsfonts}
\usepackage{graphicx}
\usepackage{wrapfig}
\usepackage{float}
\usepackage{amssymb}
\usepackage{amsmath}
\usepackage{gensymb}
\usepackage{changepage}
\usepackage[nottoc]{tocbibind}
\usepackage{tikz} 
\tikzset{font=\fontfamily{phv}\selectfont}
\usepackage{pgfplots}
\usepackage{ragged2e}
\usepackage{caption}       
\usepackage[justification=justified,singlelinecheck=false]{caption}
\DeclareCaptionLabelSeparator{mysep}{.\hspace{2mm}}
\captionsetup{format=hang, labelsep=mysep}
\usepackage{tocloft}
\usepackage{cleveref}
\usepackage{tocvsec2}
\usepackage{appendix}
\usepackage{rotating}       
\usepackage{mdframed}
\usepackage{pdfpages}      
\usepackage{multirow}      
\usepackage{booktabs}      
\usepackage{array}         
\usepackage{xr}             
\usepackage{titlesec}     


\usepackage{tocloft}        
\usepackage{nomencl}        
\makenomenclature           
\renewcommand{\nomname}{\makebox[\linewidth]{\Large{Lista de siglas, abreviaturas y símbolos \hfill}}}


\renewcommand{\cfttoctitlefont}{\raggedright\Large\bfseries} 
    \renewcommand{\cftaftertoctitle}{}
\renewcommand{\cftloftitlefont}{\raggedright\Large\bfseries}
    \renewcommand{\cftafterloftitle}{}
\renewcommand{\cftlottitlefont}{\raggedright\Large\bfseries} 
    \renewcommand{\cftafterlottitle}{}

\renewcaptionname{spanish}{\listfigurename}{Lista de figuras}
\renewcaptionname{spanish}{\listtablename}{Lista de tablas}

\setcounter{secnumdepth}{3}  
\setcounter{tocdepth}{3}    
\AtBeginDocument{%
\addtocontents{toc}{~\hfill Pág.\par}    
\addtocontents{lof}{~\hfill Pág.\par}    
\addtocontents{lot}{~\hfill Pág.\par}     
\addtocontents{apc}{~\hfill Pág.\par}
}
\counterwithout{figure}{chapter}
\counterwithout{table}{chapter}                  
\counterwithout{equation}{chapter}              


\usepackage[fitting]{tcolorbox} 
\newtcboxfit{\mybox}[1][]{%
  fit algorithm=hybrid*,
  width=12.5cm,
  height=7.5cm,
  sharp corners,
  colframe=black,colback=white,
  size=fbox,#1
}
\usepackage[scaled]{helvet}
\usepackage[T1]{fontenc}
\renewcommand\familydefault{\sfdefault}
\usepackage[eulergreek]{sansmath}
\usepackage{tgbonum}
\usepackage{setspace}
\usepackage[version=3]{mhchem}
\usepackage{textgreek}
\usepackage{geometry}
\usepackage{sectsty}
\usepackage{fancyhdr}
\usepackage{ifthen}
\setlength{\headheight}{15pt}

% Setup for headers and footers (fancyhdr)
\fancyhf{}
\lhead{\ifthenelse{\isodd{\value{page}}}{Tesis de Maestría}{Nombre del autor}}
%\rhead{Nombre del autor}
\rfoot{\thepage}
\pagestyle{fancy}



%\fancypagestyle{mypagestyle}{%
 % \fancyhf{}% Clear header/footer
 % \fancyhead[L]{Tesis de Maestría}% Author on Odd page, Centred
 % \fancyhead[L]{Nombre del autor}% Title on Even page, Centred

  
  
 % \fancyfoot[R]{\thepage}%
%\renewcommand{\baselinestretch}{1.5} 
%}



%\pagestyle{mypagestyle}

\geometry{a4paper, left=40mm, right=25mm,top=30mm,bottom=25mm,headsep=20mm}

\allsectionsfont{\normalsize}

\chapternumberfont{\raggedleft \huge}

\titleformat{\chapter}[display]
  {\normalfont\huge\bfseries}{}
  {15pt}{\huge\raggedleft}
%\patchcmd{\chapter}{\thispagestyle{plain}}{\thispagestyle{fancy}}{}{}


\makeatletter


%---------------------------
%DESCOMENTAR TODA ESTA SECCIÓN PARA QUE LOS TÍTULOS DE LOS CAPÍTULOS APAREZCAN ENUMERADOS

%\renewcommand\@chapapp{\chaptername}
%\def\@makechapterhead#1{
%  \vspace*{50\p@}
%  {\parindent  \bfseries \z@  
%    \ifnum \c@secnumdepth >\m@ne
%        {\large \centering \@chapapp\space \thechapter\par\nobreak}
%        \vskip 20\p@
%    \fi
%    \interlinepenalty\@M
%    {\raggedright \normalsize \thechapter. \space #1\par\nobreak}
%    \vskip 40\p@
%  }}
%---------------------------

\usepackage{tocloft,titlesec}
\newcommand{\listappendixname}{\hfill \large{ÍNDICE DE APÉNDICES} \hfill}
\newlistof{appendix}{apc}{\listappendixname}

\makeatletter

\AtBeginEnvironment{appendices}{%
  \clearpage
  \write\@auxout{%
    \string\let\string\latex@tf@toc\string\tf@toc
    \string\let\string\tf@toc\string\tf@apc% 
  }}
\AtEndEnvironment{appendices}{%
  \write\@auxout{%
    \string\let\string\tf@toc\string\latex@tf@toc% 
  }}
\makeatother
\usepackage{lipsum}
\pgfplotsset{compat=1.17}

\usepackage{apacite}
\usepackage{natbib}




\begin{document}
\renewcommand{\BOthers}[1]{et al.\hbox{}} %para que aparezca et al en vez de cols.


\thispagestyle{empty}
\pagenumbering{gobble}
\setstretch{1.0}
\begin{center}



\Large{\textbf{UNIVERSIDAD NACIONAL DE ASUNCIÓN\\
FACULTAD DE CIENCIAS QUÍMICAS}}
\vspace{1.5cm}

\includegraphics[width=4.8cm]{Figures/logofcqhighquality.png}\\

\vspace{1.5cm}

\textbf{\LARGE{TESIS DE MAESTRÍA}}

\vspace{1.5cm}

\textbf{\LARGE{TÍTULO DEL TRABAJO}}

\vspace{2cm}
\Large{\textbf{NOMBRE DEL POSTULANTE}}
\vspace{2cm}


\normalsize{Tesis presentada a la Dirección de Postgrado de la Facultad de Ciencias Químicas, Universidad Nacional de Asunción, como requisito para la obtención del Grado de Magíster en Ciencias de los Alimentos}

\vspace{2cm}

\textbf{\normalsize{San Lorenzo - Paraguay}}

\textbf{\normalsize{Mes - Año}}

\pagebreak
\thispagestyle{empty}

\LARGE{\textbf{UNIVERSIDAD NACIONAL DE ASUNCIÓN\\
FACULTAD DE CIENCIAS QUÍMICAS}}
\vspace{1.5cm}

\includegraphics[width=4.8cm]{Figures/logofcqhighquality.png}\\

\vspace{1.5cm}

\textbf{\LARGE{TESIS DE MAESTRÍA}}

\vspace{2cm}

\textbf{\LARGE{TÍTULO DEL TRABAJO}}

\vspace{2cm}
\Large{\textbf{MAESTRÍA EN CIENCIAS DE LOS ALIMENTOS}}

\LARGE{\textbf{NOMBRE DEL POSTULANTE
}}

\vspace{2cm}

\Large{\textbf{TUTOR/A: }}

\vspace{2cm}

\textbf{\normalsize{San Lorenzo - Paraguay}}

\textbf{\normalsize{Mes - Año}}

\end{center}


\newpage
%\pagenumbering{Roman}
%\setcounter{page}{3}
\thispagestyle{empty}

\textcolor{white}{ }

\vfill

\vspace{5cm}

\mybox{
  \vspace{0.8cm}
  \large{Autor 1}
  
     \hspace{0.4cm} \textbf{Título del trabajo.. /} Autor 1, Autor 2  - - San Lorenzo: Los Autores, 2021.
     
    \hspace{0.4cm} XX, Número de páginas p.: tabl.; 30 cm.

    \hspace{0.4cm} Trabajo de grado presentado a la Facultad de Ciencias Químicas - UNA para optar al título de CARRERA.

    \hspace{0.4cm} \textbf{1.	PALABRA CLAVE 2.  PALABRA CLAVE} I.  AUTOR.  II. Título.
                                                                                                       
\hfill \textbf{CDD: XXX}

\hfill \textbf{XXX}\\

}







\newpage
\thispagestyle{plain}


%\begin{table}[H]
%    \centering
%    \begin{tabular}{cc}
%\multirow{7}{*}{\includegraphics[width=3cm]{Figures/logofcqhighquality.png}} & \\
%& \\
%                             & \textbf{\Large{Universidad Nacional de Asunción}} \\
%                             & \\
%                             & \textbf{\Large{Facultad de Ciencias Químicas}}\\
%                             & \\
%                             & \\
%    \end{tabular}
%\end{table}

%\spacing{1.5}

%\vspace{1.5cm}
%\large
%EL PROF. \textit{Nombre del coordinador de carrera}, COORDINADOR DE TRABAJO DE GRADO DE \textit{Nombre de la carrera}, DE LA FACULTAD DE CIENCIAS QUÍMICAS DE LA UNIVERSIDAD NACIONAL DE ASUNCIÓN,

%\textbf{CERTIFICA:} que el presente documento, titulado \textbf{``TÍTULO DEL TRABAJO''}, constituye la Memoria del Trabajo de Grado que presentan los estudiantes \textbf{Autor 1} y \textbf{Autor 2} para aspirar al Título Universitario de \textit{grado de la carrera} bajo la dirección de los profesores \textbf{Orientador} y \textbf{Co-orientador}.

%Considerando que el trabajo reúne los requisitos necesarios para ser presentado ante el tribunal constituido a tal efecto y para que conste, se expide y firma el presente informe en San Lorenzo, a los x días del mes de x de 202x.


%\vfill
%\begin{center}
 %   \begin{table}[H]
  %      \centering
    %    \begin{tabular}{c}
   %     \hline\\
     %   Prof. Nombre del coordinador de carrera\\
     %   \end{tabular}
    %\end{table}
%\end{center}

\newpage
\pagenumbering{Roman}
\setcounter{page}{2}
%\thispagestyle{empty}


\large
\clearpage
\spacing{1.5}
\let\origaddvspace\addvspace
 \renewcommand{\addvspace}[1]{}
 \listoffigures
 \pagebreak
 \cleardoublepage
\addcontentsline{toc}{chapter}{Lista de siglas, abreviaturas y símbolos}
\printnomenclature[3cm]
\pagebreak 
\cleardoublepage
 \listoftables
 %\renewcommand{\addvspace}[1]{\origaddvspace{#1}}
%\pagebreak

%\cleardoublepage
%\addcontentsline{toc}{chapter}{Índice de apéndices}
%\listofappendix
\pagebreak
\cleardoublepage


\tableofcontents
%\pagebreak


\newpage
\thispagestyle{plain}
\large



\setstretch{1.5}



\raggedleft{\textbf{\huge{RESUMEN}}}



\justifying

\noindent 
\lipsum[1]


\textit{\textbf{Palabras clave:} Palabra clave 1, Palabra clave 2.}


\newpage
\thispagestyle{plain}
\large



\setstretch{1.5}



\raggedleft{\textbf{\huge{ABSTRACT}}}



\justifying

\noindent 
\lipsum[1].


\textit{\textbf{Keywords:} Palabra clave 1 en inglés, Palabra clave 2 en inglés.}


%\justify

%\large
%\clearpage
%\spacing{1.5}
%\tableofcontents
%\pagebreak

%\let\origaddvspace\addvspace
% \renewcommand{\addvspace}[1]{}
% \listoftables
% \pagebreak
% \listoffigures
% \renewcommand{\addvspace}[1]{\origaddvspace{#1}}
%\pagebreak

%\cleardoublepage
%\addcontentsline{toc}{chapter}{Índice de apéndices}
%\listofappendix
%\pagebreak
%\cleardoublepage
%\addcontentsline{toc}{chapter}{Lista de abreviaciones}
%\printnomenclature[3cm]



\pagebreak
%\lhead{Tesis de Maestría}
%\rhead{Capítulo \thechapter}
\pagenumbering{arabic}


\normalsize
\chapter[CAPÍTULO I - INTRODUCCIÓN]{INTRODUCCIÓN}


\lipsum[1]


\chapter[CAPÍTULO II - REVISIÓN DE LA LITERATURA]{REVISIÓN DE LA LITERATURA}
\citet{pulido2012análisis} \lipsum[4-13]
\begin{table}[H]
    \centering
    \caption{Caption}
    \begin{tabular}{cc}
\hline
A & B\\
\hline
       1  & 3 \\
        2 & 4\\
\hline
    \end{tabular}
     \label{tab:my_label}
\end{table}

\begin{figure}[H]
    \centering
\begin{mdframed}
\centering
    \includegraphics[width=0.1\linewidth]{Figures/logofcqhighquality.png}    
\end{mdframed}
    \caption{Caption}
    \label{fig:enter-label}
\end{figure}


\newpage
\chapter[CAPÍTULO II - OBJETIVOS]{OBJETIVOS}
\section{Objetivo general}

\lipsum[2]

\section{Objetivos específicos}
\begin{itemize}
    \item OE 1.
    \item OE 2.
    \item OE 3.
    \item OE 4.
    \item OE 5.
\end{itemize}





\newpage
\bibliographystyle{apacite}
\bibliography{Bibliografia}




%EJEMPLOS DE ABREVIACIONES
\nomenclature{cm}{Centímetros}
\nomenclature{DS}{Desviación estándar}
\nomenclature{EE}{Error estándar}




\end{document}
